% =========================================================
% PREAMBULE COMMUN — SUPPORTS DE MATHEMATIQUES
% =========================================================
% Ce fichier rassemble les packages, réglages et commandes
% partagés par le livre complet et les chapitres individuels.
% Il est chargé depuis main.tex et hérité par les sous-fichiers.

% Encodage et langue
\usepackage[utf8]{inputenc}
\usepackage[T1]{fontenc}
\usepackage{lmodern}
\usepackage[french]{babel}

% Mathématiques et symboles
\usepackage{amsmath,amssymb}
\usepackage{textcomp}
\usepackage{eurosym}

% Mise en page
\usepackage[margin=2.3cm,headheight=15pt]{geometry}
\usepackage{fancyhdr}
\usepackage{enumitem}
\usepackage{xcolor}
\usepackage[most]{tcolorbox}

% Illustrations
\usepackage{tikz}
\usetikzlibrary{arrows.meta,calc,positioning,angles,quotes}

% Liens PDF (à charger après la plupart des autres packages)
\usepackage[hidelinks]{hyperref}

% Mise en forme générale
\setlength{\parindent}{0pt}
\setlength{\parskip}{0.5em}
\setlist[itemize]{itemsep=2pt,topsep=2pt}
\setlist[enumerate]{itemsep=2pt,topsep=2pt}

% En-têtes et pieds de page
\pagestyle{fancy}
\fancyhf{}
\fancyhead[L]{\small Supports de mathématiques}
\fancyhead[R]{\small\nouppercase{\leftmark}}
\fancyfoot[C]{\small\thepage}
\renewcommand{\headrulewidth}{0.4pt}
\renewcommand{\chaptermark}[1]{\markboth{#1}{}}
\setlength{\headheight}{15pt}

% Environnements pédagogiques
\newtcolorbox{definition}{
  breakable,
  colback=blue!4,
  colframe=blue!60!black,
  fonttitle=\bfseries,
  title={Définition}
}

\newtcolorbox{propriete}[1][Propriétés]{
  breakable,
  colback=green!5!white,
  colframe=green!45!black,
  fonttitle=\bfseries,
  title={#1}
}

\newtcolorbox{attention}{
  breakable,
  colback=red!4,
  colframe=red!60!black,
  fonttitle=\bfseries,
  title={Attention}
}

\newtcolorbox{methode}{
  breakable,
  colback=orange!6,
  colframe=orange!70!black,
  fonttitle=\bfseries,
  title={Méthode}
}

% Virgule décimale (commande déjà utilisée dans certains supports)
\newcommand{\dec}{{,}}

% Mention légale : appelée dans le livre ou dans chaque chapitre autonome.
\newcommand{\MentionLegale}{%
  \begin{tcolorbox}[
    breakable,
    colback=gray!5,
    colframe=gray!55!black,
    title={Propriété intellectuelle},
    fonttitle=\bfseries
  ]
  Ce document est rédigé et mis à disposition exclusivement pour les élèves
  inscrits aux cours de Gauthier Baudelet. Toute reproduction, mise en ligne
  (Google Drive partagé, réseaux sociaux, plateformes publiques) ou transmission
  à des tiers, même à titre gratuit, est strictement interdite.
  \end{tcolorbox}%
}

% Page de mention légale pour un chapitre compilé individuellement.
\newcommand{\PageLegaleChapitre}{%
  \begin{titlepage}
    \vspace*{\fill}
    \MentionLegale
    \vspace*{\fill}
  \end{titlepage}
}

% À utiliser dans chaque fichier de chapitre avec subfiles :
% cette condition affiche la mention légale uniquement lorsque
% le chapitre est compilé seul, et pas lorsqu'il est inclus dans main.tex.